\section{Aprendizaje continuo y Memory}

\subsection{Los retos de ingeniería del Continual Learning}

\textbf{Xie Tianbao (谢天宝)}: Tomando a Cursor como ejemplo, hace un Continual Learning simple (finalización de Type), pero el ciclo no es en tiempo real (se actualiza cada varias horas), y enfrenta problemas como messy signal, gradientes efectivos con un batch size pequeño, olvido catastrófico y ataques maliciosos. El problema central es: \textbf{antes, una vez entrenado el modelo lo dejábamos freeze para hacer inference, y nunca actualizábamos los parámetros durante la inference} —esta es una diferencia de paradigma enorme, y también una línea de investigación muy grande.

\textbf{Cao Yu (曹玉)}: Desde el punto de vista lógico, ``FLOPs for Continual Learning'' es natural y coherente —cambiar FLOPs por una mejora continua de la inteligencia. Pero el reto de ingeniería es enorme: ¿con qué método hacer el Learn (¿LoRA? ¿Sparse Update?), ¿se puede lograr un Serving Per User Per Model? Espero que en 2026 haya oportunidad de conocer qué arquitecturas de modelo o qué maneras pueden realmente lograr el Test-Time Training, y no quedarse solo en la deducción teórica y en pequeños toy example.

\subsection{Dos caminos: Parameter-based vs Context-based}

\begin{knowledgebox}{Xue Fuzhao (薛福昭): sin florituras, usar Context directamente}
El Lifelong Learning tiene dos caminos:
\begin{enumerate}
  \item \textbf{Parameter-based}—LoRA / MoE con Expert añadido / Embedding Bank. Riesgo: entre más grande el modelo, más inestable; el backpropagation con un batch size pequeño es muy risky, y la experiencia de uso a largo plazo puede degradarse.
  \item \textbf{Context-based}—armar un Prompt Database (parecido a Memory Retrieval), y escribir ahí el historial. La única desventaja es que el context es muy largo—\textbf{pero un context largo, en esencia, sigue siendo FLOPs}, lo que nos regresa al principio inicial.
\end{enumerate}
Recomendación: sin florituras, hacerlo directamente como context + retrieval. La actualización del knowledge del modelo se deja al lanzamiento periódico de nuevas versiones (de cualquier forma, ahora sale un nuevo model cada varios meses).
\end{knowledgebox}

\textbf{Liu Ziwei (刘子伟)}: Un modelo grande podría parecerse a un OS—contiene distintas regiones, donde algunos parámetros deberían actualizarse lento (Low-frequency knowledge) y otros deberían actualizarse rápido (High-frequency knowledge). Desde hace ya 20--30 años existía el concepto de Slow Weight / Fast Weight (Hinton y otros). Además, distintos dominios tienen necesidades distintas respecto al Long-tail Knowledge, y ciertas key application necesariamente deben considerar este factor a nivel de modelo o de datos. Los datos tienen por naturaleza una distribución de cola larga—la parte de baja frecuencia no se mueve y la de alta frecuencia se actualiza de manera continua—esta dynamics merece un estudio a fondo.

\subsection{Resumen de la sección}
El Continual Learning es una capacidad clave para la próxima generación de IA, pero enfrenta un triple reto de estabilidad, complejidad de ingeniería y costo de Serving. El esquema Context-based es más seguro pero necesita soporte de context largo; el esquema Parameter-based es más flexible pero con mayor riesgo.

%% ===== Panorama 2026 =====
